\chapter{EXPECTED OUTCOMES}
The system is made to monitor the nutrient solution pH with a set interval of 10 minutes and regulate it within a tolerance band of ±0.2 units around the optimal setpoint defined for each crop. This is important as nutrient availability and the efficiency of root absorption are highly sensitive to pH fluctuations. By using real-time sensor feedback and automated dosing of acid solution, the system minimizes deviations caused by factors such as evaporation and nutrition dosing. \\[10pt]
The control system aims to regulate the concentration of nutrients by maintaining the EC within ±64 ppm of the crop-specific target. EC is a critical indicator of dissolved nutrient levels, and deviated values of EC can result in nutrient deficiency or overshoot. Through EC monitoring and controlled nutrients dosing, the system compensates for dilution, plant uptake, and evaporation effects. This tight control improves nutrient use efficiency, prevents excess fertilization, and ensures that plants receive a balanced nutrient supply aligned with their growth stage.\\[10pt]
The suggested technique minimizes nutrient waste that is directly connected to human dosing and recurring adjustments by automating pH and EC regulation. In order to prevent deficiencies, manual systems frequently rely on cautious over-application, which results in excess nutrient outflow during reservoir replacement. The automated method minimizes nutrition loss by applying remedial measures only when required and in exact amounts. It is anticipated that this improved nutrient management will eventually result in a 20–30\% decrease in nutrient use, which will lower operating expenses and have a less negative environmental impact.\\[10pt]
The system is designed for rapid corrective response when pH or EC values drop outside the acceptable range. When a deviation is detected, the control system starts dosing immediately to bring the parameters back within limits in under 10 minutes. This quick response time is essential to prevent long exposure of plant roots to non-optimal conditions, which can affect the growth of the plants. Rapid stabilization also demonstrates the reliability of the control strategy under various operating conditions.\\[10pt]
A key outcome of the automated system is to match or have higher fresh biomass yield reported for conventional manual hydroponic systems. By maintaining stable pH and EC levels, the system creates an optimal environment that helps in efficient nutrient absorption and healthy plant growth. Reduced stress, consistent availability of nutrients, and minimized human error, all contribute to improved growth uniformity and productivity. Achieving equal or higher yields validates the effectiveness of automation not only in resource efficiency but also in enhancing overall crop production.